% ch11_kan_extensions.tex

\chapter{Kan Extensions}

\begin{overviewbox}
You have learned how to move objects between categories (functors), how to compare
functors (natural transformations), how to build objects universally (limits and
colimits), and how to relate whole categories (adjunctions). There is one
construction that quietly underlies all of these: the \textbf{Kan extension}.

When a functor $F : \mathcal{C} \to \mathcal{D}$ cannot be defined on a larger
category $\mathcal{E}$ without additional choices, a Kan extension asks: what is
the \emph{best possible approximation} to extending $F$ along a functor
$p : \mathcal{C} \to \mathcal{E}$? The answer turns out to subsume limits,
colimits, and adjunctions as special cases. Saunders Mac Lane wrote that Kan
extensions are ``perhaps the most important concept of category theory.''

This chapter builds Kan extensions from that intuition outward: we start with
the idea of approximation, develop the colimit formula for left Kan extensions,
the limit formula for right Kan extensions, and then reveal how Kan extensions
themselves are adjoint functors between functor categories.
\end{overviewbox}

% ─────────────────────────────────────────────────────────────────────────────
\begin{nameorigin}
\term{Kan extension} is named for \emph{Daniel M.\ Kan} (1927--2013,
Israeli-American mathematician at MIT). Kan introduced the construction
in his 1958 paper ``Adjoint Functors,'' as part of the same project that
gave us adjunctions. The original motivation: extending a functor
$F : \mathcal{C} \to \mathcal{D}$ along a functor $p : \mathcal{C} \to
\mathcal{E}$ to a functor $G : \mathcal{E} \to \mathcal{D}$ in the
``best'' way (left or right universal).

Kan extensions subsume limits, colimits, and adjunctions as special
cases (Mac Lane's famous slogan: ``Kan extensions are perhaps the most
important concept of category theory''). Kan's other contributions to
the foundations of homotopy theory and category theory are extensive:
Kan complexes (Vol IV Ch.~39), the Kan model structure on simplicial
sets (Vol IV Ch.~41), and the now-standard categorical conception of
adjoint functors.
\end{nameorigin}

\section{Extending Along a Functor}

\subsection{Informally}

Suppose you have a small category $\mathcal{C}$, a functor
$p : \mathcal{C} \to \mathcal{E}$, and a functor $F : \mathcal{C} \to \mathcal{D}$.
You would like a functor $G : \mathcal{E} \to \mathcal{D}$ that ``agrees with $F$
on $\mathcal{C}$,'' meaning $G \circ p \cong F$.

\begin{center}
\begin{tikzcd}
\mathcal{C}
  \arrow[r, "p"]
  \arrow[dr, "F"']
& \mathcal{E}
  \arrow[d, dashed, "G"]
\\
& \mathcal{D}
\end{tikzcd}
\end{center}

Most of the time there is no \emph{unique} such $G$. Many functors could fill the
dashed arrow. So we ask a softer question: among all functors
$G : \mathcal{E} \to \mathcal{D}$ equipped with a natural transformation
$\alpha : F \Rightarrow G \circ p$, is there a \emph{universal} one — one that every
other such pair factors through uniquely?

That universal pair is called the \term{left Kan extension} of $F$ along $p$,
written $\operatorname{Lan}_p F$. Its universal property says: any natural
transformation $F \Rightarrow G \circ p$ factors uniquely through the canonical
transformation $\eta : F \Rightarrow (\operatorname{Lan}_p F) \circ p$.

Dually, the \term{right Kan extension} $\operatorname{Ran}_p F$ is equipped with
a natural transformation $\varepsilon : (\operatorname{Ran}_p F) \circ p \Rightarrow F$,
and any transformation $G \circ p \Rightarrow F$ factors uniquely through $\varepsilon$.

\begin{dialog}
``Is this just like an adjunction?''

Exactly. The Kan extension is defined by a universal property, and universal
properties always define adjoints in disguise. We will make that precise in
Section~\ref{sec:kan-as-adjoint}.
\end{dialog}

\subsection{Expanding}

Let us think about what ``best approximation'' means concretely. The functor $p$
maps objects of $\mathcal{C}$ into $\mathcal{E}$. If $e \in \mathcal{E}$ happens
to equal $p(c)$ for some $c$, then $G(e)$ had better relate to $F(c)$. But what
if $e$ is not in the image of $p$ at all?

A natural answer is: look at all the ways $e$ can be ``approached'' by objects
in the image of $p$. Form a diagram indexed by those approaches, evaluate $F$ on
the objects of $\mathcal{C}$ involved, and take the colimit of that diagram.
The colimit assembles all the relevant information about $e$ from $F$'s perspective.

This is the \textbf{colimit formula}:
\[
  (\operatorname{Lan}_p F)(e) \;=\; \operatorname{colim}_{(c,\, p(c) \to e)\,\in\, (p \downarrow e)} F(c).
\]
The comma category $(p \downarrow e)$ has objects $(c, \phi)$ where
$c \in \mathcal{C}$ and $\phi : p(c) \to e$ is a morphism in $\mathcal{E}$.
It collects all the ``windows into $e$'' through the image of $p$.

Dually, the right Kan extension uses a limit:
\[
  (\operatorname{Ran}_p F)(e) \;=\; \lim_{(c,\, e \to p(c))\,\in\, (e \downarrow p)} F(c).
\]

\begin{dialog}
``Why colimit for left and limit for right?''

Left Kan extensions are left adjoints, and left adjoints preserve colimits — so
the formula is built from colimits. Right Kan extensions are right adjoints, and
right adjoints preserve limits — so the formula uses limits. The direction of the
arrow in the comma category tracks which side you are on.
\end{dialog}

\subsection{Formalizing}

\begin{definition}
Let $p : \mathcal{C} \to \mathcal{E}$ and $F : \mathcal{C} \to \mathcal{D}$ be
functors. A \term{left Kan extension} of $F$ along $p$ is a functor
$\operatorname{Lan}_p F : \mathcal{E} \to \mathcal{D}$ together with a natural
transformation
\[
  \eta : F \;\Rightarrow\; (\operatorname{Lan}_p F) \circ p,
\]
called the \term{unit}, satisfying the following universal property: for every
functor $G : \mathcal{E} \to \mathcal{D}$ and every natural transformation
$\alpha : F \Rightarrow G \circ p$, there exists a unique natural transformation
$\bar{\alpha} : \operatorname{Lan}_p F \Rightarrow G$ such that
$(\bar{\alpha} \circ p) \cdot \eta = \alpha$.

\begin{howtosay}
$\operatorname{Lan}_p F$ is read ``the left Kan extension of $F$ along $p$.''
The subscript $p$ names the functor you are extending along.
\end{howtosay}
\end{definition}

\begin{definition}
A \term{right Kan extension} of $F$ along $p$ is a functor
$\operatorname{Ran}_p F : \mathcal{E} \to \mathcal{D}$ together with a natural
transformation
\[
  \varepsilon : (\operatorname{Ran}_p F) \circ p \;\Rightarrow\; F,
\]
called the \term{counit}, satisfying the dual universal property: for every functor
$G : \mathcal{E} \to \mathcal{D}$ and every natural transformation
$\beta : G \circ p \Rightarrow F$, there exists a unique natural transformation
$\bar{\beta} : G \Rightarrow \operatorname{Ran}_p F$ such that
$\varepsilon \cdot (\bar{\beta} \circ p) = \beta$.
\end{definition}

% ─────────────────────────────────────────────────────────────────────────────
\section{The Colimit Formula}

\subsection{Formalizing}

The existence of left Kan extensions, and an explicit formula for computing them,
rests on the colimit formula stated informally above. We now make it precise.

\begin{definition}
The \term{comma category} $(p \downarrow e)$, for a functor
$p : \mathcal{C} \to \mathcal{E}$ and an object $e \in \mathcal{E}$, has:
\begin{itemize}
  \item \textbf{Objects:} pairs $(c, \phi)$ where $c \in \mathcal{C}$ and
    $\phi : p(c) \to e$ in $\mathcal{E}$.
  \item \textbf{Morphisms:} $(c, \phi) \to (c', \phi')$ is a morphism
    $f : c \to c'$ in $\mathcal{C}$ such that $\phi' \circ p(f) = \phi$.
\end{itemize}
There is a forgetful functor $\pi : (p \downarrow e) \to \mathcal{C}$
sending $(c, \phi) \mapsto c$.
\end{definition}

\begin{theorem}[Colimit Formula for Left Kan Extensions]
If $\mathcal{D}$ has all small colimits and $\mathcal{C}$ is small, then
$\operatorname{Lan}_p F$ exists and is given pointwise by
\[
  (\operatorname{Lan}_p F)(e) \;\cong\; \operatorname{colim}\bigl(F \circ \pi_{e}\bigr),
\]
where $\pi_e : (p \downarrow e) \to \mathcal{C}$ is the forgetful functor.
\end{theorem}

The colimit is over a diagram that assembles all $F(c)$ for all morphisms
$p(c) \to e$. As $e$ varies, these diagrams fit together naturally to make
$\operatorname{Lan}_p F$ a functor.

\begin{theorem}[Limit Formula for Right Kan Extensions]
Dually, if $\mathcal{D}$ has all small limits, then $\operatorname{Ran}_p F$ exists
and is given pointwise by
\[
  (\operatorname{Ran}_p F)(e) \;\cong\; \lim\bigl(F \circ \pi^e\bigr),
\]
where $\pi^e : (e \downarrow p) \to \mathcal{C}$ is the forgetful functor from the
comma category $(e \downarrow p)$ whose objects are $(c, \phi : e \to p(c))$.
\end{theorem}

\subsection{Reorganizing: Limits and Colimits Are Kan Extensions}

The colimit formula makes a striking observation possible. Take $\mathcal{E} = \mathbf{1}$,
the category with one object and one morphism, and let $p : \mathcal{C} \to \mathbf{1}$
be the unique functor. Then:
\begin{itemize}
  \item The comma category $(p \downarrow \star)$, for the single object $\star \in \mathbf{1}$,
    is isomorphic to $\mathcal{C}$ itself (every morphism $p(c) \to \star$ is the
    identity).
  \item So $(\operatorname{Lan}_p F)(\star) = \operatorname{colim}_{\mathcal{C}} F$.
\end{itemize}

\begin{center}
\begin{tikzcd}
\mathcal{C}
  \arrow[r, "p"]
  \arrow[dr, "F"']
& \mathbf{1}
  \arrow[d, dashed, "\operatorname{Lan}_p F"]
\\
& \mathcal{D}
\end{tikzcd}
\end{center}

In other words, \emph{colimits are left Kan extensions along the unique functor
to the terminal category}, and limits are right Kan extensions along the same functor.
Kan extensions are more general: they push a functor along an arbitrary functor,
not just along the collapse to a point.

% ─────────────────────────────────────────────────────────────────────────────
\section{Kan Extensions as Adjoints}
\label{sec:kan-as-adjoint}

\subsection{Formalizing}

Fix categories $\mathcal{C}$, $\mathcal{D}$, $\mathcal{E}$ and a functor
$p : \mathcal{C} \to \mathcal{E}$. There is a \term{restriction functor}
\[
  (-) \circ p \;:\; [\mathcal{E}, \mathcal{D}] \;\longrightarrow\; [\mathcal{C}, \mathcal{D}]
\]
that sends a functor $G : \mathcal{E} \to \mathcal{D}$ to the composite
$G \circ p : \mathcal{C} \to \mathcal{D}$.

\begin{theorem}
When Kan extensions exist, they form adjoint pairs with the restriction functor:
\begin{align*}
  \operatorname{Lan}_p &\dashv (-) \circ p, \\
  (-) \circ p &\dashv \operatorname{Ran}_p.
\end{align*}
In explicit terms, there are natural bijections:
\begin{align*}
  [\mathcal{E}, \mathcal{D}](\operatorname{Lan}_p F,\, G)
    &\;\cong\; [\mathcal{C}, \mathcal{D}](F,\, G \circ p), \\
  [\mathcal{C}, \mathcal{D}](G \circ p,\, F)
    &\;\cong\; [\mathcal{E}, \mathcal{D}](G,\, \operatorname{Ran}_p F).
\end{align*}
\end{theorem}

\begin{proof}[Proof sketch]
The hom-set bijection for $\operatorname{Lan}_p \dashv (-) \circ p$ is precisely
the universal property of the left Kan extension: natural transformations
$\operatorname{Lan}_p F \Rightarrow G$ are in natural bijection with natural
transformations $F \Rightarrow G \circ p$, mediated by pre- or post-composition
with the unit $\eta$.
\end{proof}

This is the key insight: \emph{forming a left Kan extension is left adjoint to
restricting along $p$}. The construction that was built by universal property
is, in the language of the previous chapters, exactly an adjunction between
functor categories.

\subsection{Simplifying}

The four main constructions of the book now appear as Kan extensions:

\begin{center}
\renewcommand{\arraystretch}{1.4}
\begin{tabular}{lll}
\hline
\textbf{Construction} & \textbf{As a Kan extension} & \textbf{Along} \\
\hline
Colimit of $F$ & $\operatorname{Lan}_p F$  & $p : \mathcal{C} \to \mathbf{1}$ \\
Limit of $F$   & $\operatorname{Ran}_p F$  & $p : \mathcal{C} \to \mathbf{1}$ \\
Left adjoint of $G$ & $\operatorname{Lan}_Y \operatorname{id}$ & Yoneda embedding $Y$ \\
Right adjoint of $G$ & $\operatorname{Ran}_Y \operatorname{id}$ & Yoneda embedding $Y$ \\
\hline
\end{tabular}
\end{center}

% ─────────────────────────────────────────────────────────────────────────────
\section{Density and the Density Corollary}

\subsection{Formalizing}

A functor $p : \mathcal{C} \to \mathcal{E}$ is called \term{dense} if every
object of $\mathcal{E}$ is a colimit of objects in the image of $p$. The precise
statement involves Kan extensions:

\begin{definition}
A functor $p : \mathcal{C} \to \mathcal{E}$ is \term{dense} if the left Kan
extension of $p$ along itself is naturally isomorphic to the identity:
\[
  \operatorname{Lan}_p p \;\cong\; \operatorname{Id}_{\mathcal{E}}.
\]
\end{definition}

The Yoneda embedding $Y : \mathcal{C} \to [\mathcal{C}^{op}, \mathbf{Set}]$ is
dense. This gives the \textbf{density corollary}: every presheaf is a canonical
colimit of representable presheaves.

\begin{theorem}[Density of the Yoneda Embedding]
For any small category $\mathcal{C}$ and any presheaf $F \in [\mathcal{C}^{op}, \mathbf{Set}]$:
\[
  F \;\cong\; \operatorname{colim}_{(c,\, Y(c) \to F)\,\in\, (Y \downarrow F)} Y(c).
\]
\end{theorem}

In words: every presheaf $F$ is the colimit of all representable presheaves
mapping into it. The ``shape'' of a presheaf is entirely determined by how
representables map into it — which is the content of the Yoneda lemma stated
from a new angle.

\begin{dialog}
``So the Yoneda lemma and Kan extensions are the same thing?''

Not quite, but they are deeply connected. The Yoneda lemma says that
$[\mathcal{C}^{op}, \mathbf{Set}](Y(c), F) \cong F(c)$. The density formula
says that $F$ is the left Kan extension of itself along $Y$. Together, they
say that presheaf categories are the canonical home where every object is
freely built from representables.
\end{dialog}

\subsection{Formally}

\begin{theorem}[Every Functor is a Left Kan Extension]
For any functor $F : \mathcal{C} \to \mathcal{D}$:
\[
  F \;\cong\; \operatorname{Lan}_{\operatorname{Id}_\mathcal{C}} F.
\]
That is, $F$ is its own left Kan extension along the identity. More substantively,
if $Y : \mathcal{C} \to \hat{\mathcal{C}}$ is the Yoneda embedding, then:
\[
  \operatorname{Lan}_Y F \;\cong\; F^+,
\]
where $F^+ : \hat{\mathcal{C}} \to \mathcal{D}$ is the unique (up to isomorphism)
colimit-preserving extension of $F$ to the presheaf category.
\end{theorem}

This is the universal property of the presheaf category: $\hat{\mathcal{C}}$ is
the free cocompletion of $\mathcal{C}$.

% ─────────────────────────────────────────────────────────────────────────────
\section{Exercises}

\begin{exercisesbox}
\begin{qa}
\qaitem{What is the left Kan extension $\operatorname{Lan}_p F$ of $F$ along $p$?}
       {A functor $\operatorname{Lan}_p F : \mathcal{E} \to \mathcal{D}$ with a unit
        $\eta : F \Rightarrow (\operatorname{Lan}_p F) \circ p$, universal among all
        such pairs.}

\qaitem{What is the universal property of the unit $\eta$?}
       {Any $\alpha : F \Rightarrow G \circ p$ factors uniquely through $\eta$
        via a natural transformation $\bar{\alpha} : \operatorname{Lan}_p F \Rightarrow G$.}

\qaitem{How do you read $\operatorname{Lan}_p F$?}
       {``The left Kan extension of $F$ along $p$.'' The subscript names the
        functor you are extending along.}

\qaitem{What is the colimit formula for $(\operatorname{Lan}_p F)(e)$?}
       {$\operatorname{colim}\bigl(F \circ \pi_e\bigr)$, where $\pi_e : (p \downarrow e) \to \mathcal{C}$
        is the forgetful functor from the comma category.}

\qaitem{What are the objects of the comma category $(p \downarrow e)$?}
       {Pairs $(c, \phi)$ where $c \in \mathcal{C}$ and $\phi : p(c) \to e$ in $\mathcal{E}$.}

\qaitem{What are the morphisms in $(p \downarrow e)$?}
       {A morphism $(c, \phi) \to (c', \phi')$ is $f : c \to c'$ in $\mathcal{C}$
        such that $\phi' \circ p(f) = \phi$.}

\qaitem{What is the limit formula for $(\operatorname{Ran}_p F)(e)$?}
       {$\lim\bigl(F \circ \pi^e\bigr)$, where $\pi^e : (e \downarrow p) \to \mathcal{C}$
        is the forgetful functor from the comma category whose objects have
        $\phi : e \to p(c)$.}

\qaitem{How does a colimit of $F : \mathcal{C} \to \mathcal{D}$ arise as a Kan extension?}
       {As $(\operatorname{Lan}_p F)(\star)$ where $p : \mathcal{C} \to \mathbf{1}$
        is the unique functor.}

\qaitem{What is the adjunction involving $\operatorname{Lan}_p$?}
       {$\operatorname{Lan}_p \dashv (-) \circ p$, as functors between functor
        categories $[\mathcal{C}, \mathcal{D}]$ and $[\mathcal{E}, \mathcal{D}]$.}

\qaitem{What is the adjunction involving $\operatorname{Ran}_p$?}
       {$(-) \circ p \dashv \operatorname{Ran}_p$.}

\qaitem{What does it mean for a functor $p : \mathcal{C} \to \mathcal{E}$ to be dense?}
       {$\operatorname{Lan}_p p \cong \operatorname{Id}_{\mathcal{E}}$; every object
        of $\mathcal{E}$ is a canonical colimit of objects in the image of $p$.}

\qaitem{Is the Yoneda embedding $Y : \mathcal{C} \to \hat{\mathcal{C}}$ dense?}
       {Yes. Every presheaf is a colimit of representable presheaves.}

\qaitem{What is the density corollary?}
       {Every presheaf $F$ is the colimit of all representables $Y(c)$ over
        the comma category $(Y \downarrow F)$.}

\qaitem{What is the free cocompletion of $\mathcal{C}$?}
       {The presheaf category $\hat{\mathcal{C}} = [\mathcal{C}^{op}, \mathbf{Set}]$:
        the universal category that receives $\mathcal{C}$ and has all colimits.}

\qaitem{If $\mathcal{D}$ has all colimits, what is the unique colimit-preserving
        extension of $F : \mathcal{C} \to \mathcal{D}$ to $\hat{\mathcal{C}}$?}
       {$\operatorname{Lan}_Y F$, the left Kan extension of $F$ along the Yoneda embedding.}

\qaitem{What does $\operatorname{Lan}_p F$ compute at $e$ when every object of
        $\mathcal{E}$ is in the image of $p$?}
       {It essentially evaluates $F$ at a preimage of $e$; when $p$ is an
        isomorphism the formula recovers $F \circ p^{-1}$.}

\qaitem{Can $\operatorname{Lan}_p F$ fail to exist?}
       {Yes, when $\mathcal{D}$ lacks the required colimits. The colimit formula
        requires $\mathcal{D}$ to have all small colimits.}

\qaitem{What is the restriction functor?}
       {$(-) \circ p : [\mathcal{E}, \mathcal{D}] \to [\mathcal{C}, \mathcal{D}]$,
        precomposing every functor $\mathcal{E} \to \mathcal{D}$ with $p$.}

\qaitem{Which Kan extension is left adjoint to restriction, and which is right?}
       {$\operatorname{Lan}_p$ is left adjoint; $\operatorname{Ran}_p$ is right adjoint.}

\qaitem{What does it mean that ``every concept is a Kan extension''?}
       {Limits, colimits, adjoints, and many other universal constructions can each
        be expressed as a Kan extension of an appropriate functor along an appropriate
        functor.}

\qaitem{How does the right Kan extension $\operatorname{Ran}_p F$ relate to adjoints?}
       {If $F$ has a right adjoint $G$, and $p = F$, then $\operatorname{Ran}_F \operatorname{Id}
        \cong G$; right adjoints can be computed as right Kan extensions along the
        left adjoint.}

\qaitem{What structure does the comma category $(p \downarrow e)$ have as $e$ varies?}
       {It varies functorially: a morphism $e \to e'$ in $\mathcal{E}$ induces a
        functor $(p \downarrow e) \to (p \downarrow e')$ by composing with the morphism.}

\qaitem{What is the unit $\eta$ for a left Kan extension?}
       {A natural transformation $\eta : F \Rightarrow (\operatorname{Lan}_p F) \circ p$;
        it is the canonical ``comparison'' map from $F$ to its extension composed with $p$.}

\qaitem{What is the counit $\varepsilon$ for a right Kan extension?}
       {A natural transformation $\varepsilon : (\operatorname{Ran}_p F) \circ p \Rightarrow F$;
        the canonical comparison in the other direction.}

\qaitem{If $p$ is an equivalence of categories, what can you say about
        $\operatorname{Lan}_p F$?}
       {It is isomorphic to $F \circ p^{-1}$ (the composite with a quasi-inverse),
        and the unit $\eta$ is a natural isomorphism.}

\qaitem{What condition on $p$ makes every functor $F$ compute its Kan extension
        without taking genuine colimits?}
       {If $p$ is fully faithful, then $(\operatorname{Lan}_p F)(p(c)) \cong F(c)$:
        the extension recovers $F$ exactly on the image of $p$.}

\qaitem{What is a pointwise Kan extension?}
       {A Kan extension computed by the colimit (or limit) formula at each object
        separately. Most Kan extensions in practice are pointwise.}

\qaitem{Why are pointwise Kan extensions better-behaved?}
       {They are preserved by representable functors; equivalently, they interact
        well with limits and colimits in $\mathcal{D}$.}

\qaitem{How does sheafification relate to Kan extensions?}
       {Sheafification is the left Kan extension of the inclusion of sheaves into
        presheaves, along the identity on presheaves — equivalently, the left adjoint
        to the forgetful functor from sheaves to presheaves.}

\qaitem{What single adjunction encodes both $\operatorname{Lan}_p$ and $\operatorname{Ran}_p$?}
       {There is no single adjunction; they form two separate adjunctions:
        $\operatorname{Lan}_p \dashv (-)\circ p \dashv \operatorname{Ran}_p$,
        making $(-) \circ p$ a functor with a left and a right adjoint.}
\end{qa}
\end{exercisesbox}

% ─────────────────────────────────────────────────────────────────────────────
\begin{summarybox}
A \textbf{Kan extension} of $F : \mathcal{C} \to \mathcal{D}$ along
$p : \mathcal{C} \to \mathcal{E}$ is the best approximation to extending $F$
to $\mathcal{E}$. The \textbf{left Kan extension} $\operatorname{Lan}_p F$ is
computed by a colimit over the comma category $(p \downarrow e)$ at each
object $e \in \mathcal{E}$; the \textbf{right Kan extension} $\operatorname{Ran}_p F$
by a limit. Both are adjoints to the restriction functor $(-) \circ p$.

Limits and colimits are the special case where $\mathcal{E} = \mathbf{1}$.
The Yoneda embedding is dense, giving the density corollary: every presheaf
is canonically a colimit of representables. The presheaf category is the free
cocompletion of $\mathcal{C}$.
\end{summarybox}

% ─────────────────────────────────────────────────────────────────────────────
\section{Formal Restatement}

\begin{formalrestatement}
\textbf{Left Kan extension.} Given $p : \mathcal{C} \to \mathcal{E}$ and
$F : \mathcal{C} \to \mathcal{D}$, the pair $(\operatorname{Lan}_p F, \eta)$
where $\operatorname{Lan}_p F : \mathcal{E} \to \mathcal{D}$ and
$\eta : F \Rightarrow (\operatorname{Lan}_p F) \circ p$ is terminal in the category
of such pairs, yielding a natural bijection
\[
  [\mathcal{E}, \mathcal{D}](\operatorname{Lan}_p F,\, G) \;\cong\;
  [\mathcal{C}, \mathcal{D}](F,\, G \circ p).
\]

\textbf{Colimit formula.} If $\mathcal{C}$ is small and $\mathcal{D}$ is cocomplete:
\[
  (\operatorname{Lan}_p F)(e) \;=\; \operatorname{colim}\bigl((p \downarrow e)
  \xrightarrow{\pi_e} \mathcal{C} \xrightarrow{F} \mathcal{D}\bigr).
\]

\textbf{Right Kan extension.} Dually, $(\operatorname{Ran}_p F, \varepsilon)$ with
$\varepsilon : (\operatorname{Ran}_p F) \circ p \Rightarrow F$, and:
\[
  [\mathcal{C}, \mathcal{D}](G \circ p,\, F) \;\cong\;
  [\mathcal{E}, \mathcal{D}](G,\, \operatorname{Ran}_p F).
\]

\textbf{Limit formula.} If $\mathcal{D}$ is complete:
\[
  (\operatorname{Ran}_p F)(e) \;=\; \lim\bigl((e \downarrow p)
  \xrightarrow{\pi^e} \mathcal{C} \xrightarrow{F} \mathcal{D}\bigr).
\]

\textbf{Adjunction triple.} $\operatorname{Lan}_p \dashv (-) \circ p \dashv \operatorname{Ran}_p$
as functors $[\mathcal{C}, \mathcal{D}] \rightleftarrows [\mathcal{E}, \mathcal{D}]$.

\textbf{Density.} $p$ is dense iff $\operatorname{Lan}_p p \cong \operatorname{Id}_\mathcal{E}$.
The Yoneda embedding $Y : \mathcal{C} \to \hat{\mathcal{C}}$ is dense;
$\hat{\mathcal{C}}$ is the free cocompletion of $\mathcal{C}$.
\end{formalrestatement}
